\sectiontitle{Section VI:}

                                  Section VI:
                                            Selection and Stability
Ever since self-replicating molecules came about, they have been producing like mad
and proliferating in ever more varieties. Moreover, have been agglomerating in
gigantic colonies that, seen at their own level, are also self-replicating entities. We
ourselves are huge self-replicating molecule-heaps. Ever upward builds this dizzying
spire of self-replicating structures. What gives this whole movement any coherent
direction? How and why does complexity evolve from simplicity? The gist of the
answer is that any active organism engages in external behavior that has repercussions
back on the organism. These repercussions then influence the further behavior of the
organism, which in turn engenders further repercussions, and so on. Such feedback
selectively reinforces certain kinds of structures and strategies while suppressing
others. Thus through feedback, certain types of organism are stabilized and can
become the building blocks for higher-level organisms. In this hierarchical way, very
complex yet stable structures and strategies can emerge out of a very 'Primitive
bottom level. What is the nature of the structures and strategies that emerge? How
stable are they? How inevitable? How arbitrary? How do . game-theoretical models
and computer simulations shed light on the competition of many organisms in a
society? Are all such organisms at odds with each other, acting maximally selfish? Or
need selfish organisms- always be at odds with each other? Can cooperation and even
a seeming morality verge purely as a consequence of the laws that govern self-
replication and the universe's impersonal preference for various styles? The iterated
Prisoner's Dilemma, a pithy idealization of such questions, is the focus of the last of
these three chapters.




The Genetic Code: Arbitrary?                                                       670


                   The Genetic Code: Arbitrary?
                                     March, 1982

It all began with a pesky student of mine named Vahe Sarkissian. I was telling a class
about one of my favorite notions: the analogy between the complex machinery in a
living cell that enables a DNA molecule to replicate itself, and the clever machinery
in a mathematical system that enables a formula to say things about itself. To my
mind, the resemblance is deep and fruitful; it has afforded me sharper insights into
both domains. Although Vahe appreciated the analogy, he doubted the validity of one
important aspect of it, and so he brought the matter up in class. His challenge forced
me to think the issues through carefully, and en route I encountered some fascinating
details of molecular biology that I might otherwise never, have known. What I find
gratifying is how quickly people can come to appreciate these intricacies without
having studied molecular biology. I'll therefore attempt to sketch out the necessary
background, and then I'll explain the problem and what I believe to be its resolution. .
Both of the profound twentieth-century discoveries involved in this analogy depend
crucially on codes: curiously arbitrary-seeming mappings from one set of entities to
another set of entities. In metamathematics, the code is Godel numbering; in biology,
it is the genetic code. In Godel numbering, invented by Kurt Godel in about 1930,
code numbers are assigned to various mathematical symbols (plus signs, digits, and
parentheses, for example), just as license numbers are assigned to cars or telephone
area codes to cities. The symbol '0', for instance, might be associated with the quantity
666, and a formula like "0=0" might inherit' the number 666,111,666 from its
constituent symbols. Gödel’s mapping connects entities from two intrinsically
unrelated domains, one typographical (printed symbols) and the other abstract
(Platonic numbers), and allows any system that can talk about numbers to talk about
itself: in code.
         The genetic code is likewise a mapping between two mutually unrelated
domains. In this case, though, both domains consist of chemical units. To someone
unfamiliar with chemical terminology, the two domains might sound so similar that
the connection of one with the other would appear mundane. But actually, nobody had
ever in the least suspected that one set




The Genetic Code: Arbitrary?                                                         671

of chemicals could code for another set. Indeed, the very idea is somewhat baffling: If
there is a code, then who invented it? What kinds of messages are written in it? Who
writes them? Who reads them? Why not just write the messages directly, rather than
in code? Are there lots of codes? Are they arbitrary? I would hope that this list of
who's and why's would perk up anybody's curiosity, and make them see that a code
connecting two unrelated families of chemicals is not to be taken for granted.

                                       *   *    *

         Over the past few billion years, a scheme gradually evolved in living beings
according to which a unit of one chemical "species" is assigned as a code for a unit of
another "species". Actually, a triplet of units of what I will call Species I is assigned
to a unit of Species II. In fact (just to make things impossibly complex), sometimes
several different triplets are assigned! But for the moment, that is beside the point.
The main fact is simply that members of two entirely unrelated species of chemical
units are "mapped" onto each other. Vahe's question had to do with how arbitrary this
mapping really is. I claimed it was arbitrary, while he claimed there must be some
comprehensible rhyme or reason to it.
         Species I is the nucleotides. Species II is the amino acids. If these words are
not in your vocabulary, don't panic! In fact, you are very likely my ideal audience.
You need not know Word One about chemistry to be able to imagine this
correspondence, this match-up between members of two different chemical species.
All you need to know is that to each triplet of nucleotides (whatever nucleotides might
be), there is matched an amino acid (whatever that is!). That is what the genetic code
is.
         I mentioned area codes earlier. They are useful to keep in mind, since they too
involve triplets (of digits). "212" codes for New York City, "619" for San Diego, and
so on. Clearly, these connections are not intrinsic: "619" could as easily have been
New York's area code. It is hard to imagine two more disparate domains than three-
digit numbers and geographical areas. Yet this mapping-this "telephonic code"-is one"
of the most useful we know!
         The purpose of the genetic code would be hard to describe without adding a
little about the constituents of the cell. The "personality", or character, of a cell is
stored in its genes. But genes are essentially static, like words in a book. For them to
come alive and have observable effects, they must be translated into dynamic agents.
These agents are proteins, and their actions realize the potential of the genes. They
"express" the genes and thereby create the character of the cell. As it turns out, genes
are strings of nucleotides, and proteins are strings of amino acids. The cell's
personality is therefore written in the passive chemical units of Species I. Through the
genetic code, this description can be converted into a vast population of dynamic
agents made out of chemical units of Species II. Hence, thanks to




The Genetic Code: Arbitrary?                                                         672

the genetic code, the cell's personality, implicitly defined by its genes, can emerge and
bloom.
         There are twenty different amino acids, so you-might think there would be
twenty different nucleotide triplets. It's not quite that simple. There happen to be four
different nucleotides involved in the genetic code, denoted 'A', 'C', 'G', and 'U' (which
stand for adenine, cytosine, guanine, and uracil ). Every possible triplet (beginning
with AAA, AAC, AAG and going all the way to UUU) stands for some amino acid.
(Well, not quite. Three triplets do not, but for the moment that is just a detail.) How
many such triplets are there? Sixty-four, of course: 4 X 4 X 4. So 61 (64 -3) different
triplets are matched up with twenty amino acids. Consequently, most amino acids are
coded for by more than one codon (a triplet of nucleotides). This would be like having
most cities represented by several different "area codons", all just as good as each
other. Indeed, there are some amino acids that have six different codons, some that
have four, some that have three, and some that have two; only a couple have one. The
complete genetic code is shown, for your convenience, in Figure 27-1.

        FIGURE 27-1. The genetic code. A typical codon such as "CA U" is seen to
represent the amino acid histidine. Notice the redundancy and partial symmetry of
this chart. Are such features important or necessary? Could this chart have been
different but life the same?




The Genetic Code: Arbitrary?                                                         673

Back to Vahe and me. I had pointed out to my class that Gödel’s numbering scheme
was quite arbitrary. Gödel could have made just about any number correspond to each
of the mathematical symbols involved; it would not have made the slightest difference
to the success of his work. I had then said that the analogous statement would hold for
the genetic code. But Vahe had the feeling that the genetic code, so tied in with the
secret of life, must be deeper. It seemed to him intuitively that each amino acid should
be related to its particular codon (or codons) for some compelling reason-that there
must be some fundamental chemical necessity for the relation. To caricature Vahe's
position, one might say: "Gödel’s code is the work of a mere mortal, but the genetic
code is the work of God. Therefore the genetic code must be perfect, inevitable, and
unalterable."
        I was quick to retort that, as far as I could tell, such was not the case. I said
that the genetic code seemed to me every bit as arbitrary as the way the telephone
company assigned area codes, every bit as arbitrary as Gödel’s numbering scheme.
On the blackboard I drew some pictures of the molecules involved, to show my
reasons for thinking this way. But as I stood there at the board, a few things began to
nag at me, places where I was not entirely sure of what I was saying, "facts" I knew I
really ought to check up on. This new desire to prove to my class the arbitrary nature
of the genetic code then led me down some fascinating paths in molecular biology,
and my findings are what I wish to report here.

                                       *    *    *
        Let us return to the cell. A cell is a little hotbed of activity, rather like a
miniature town. There are basically two kinds of entity in this town. There are passive
objects, "lumps" that just sit around and wait for somebody to do something to them,
and there are active agents, "doers" who always want to get in there and do
something. These active agents are, for the most part, enzymes. (The very definition
of an enzyme is that it is a protein that does something.) Each enzyme has a specific
job it can carry out, and it does this job on lumps of a specific type. (Actually,
enzymes can even act on other enzymes. However, enzymes being acted on have to
act "lumpish" during the operation, like normally vigorous patients anesthetized in a
dentist's chair. Once they're released, they may well re-become active doers.) How an
enzyme works is not our business here. We can just assume that enzymes do their
thing, whether it involves splitting some lump in two, welding two lumps together,
transporting some lump from one place to another, or performing some other
chemical act.
        A marvelous thing about cells is that they are so elegantly designed that for
many purposes one can totally ignore their chemistry and think just about their logic.
In fact, that is the only way I know to think about the goings-on inside cells, since I
am not a biochemist. Although I use the chemical names for things, my true image,
deep down inside, is hardly one of chemicals. It is really one of little objects that
somehow magically behave




The Genetic Code: Arbitrary?                                                         674

in certain ways specified in biochemistry books. My view of the chemicals involved
in the processes of life is like most people's view of cars: they know what cars will do
in all kinds of situations, but they don't really understand how cars work. I get a kick
out of batting around technical terms of biochemistry when in fact I understand only
their logic-"the molecular logic of the living state", as Albert Lehninger calls it. The
fact that one can get away with this is one of the beauties of molecular biology, and it
is this beauty that we are celebrating here.

*   *   *

         An enzyme is just a kind of protein, and all proteins are strangely curled
molecules made up of amino acids. The curliness is crucial. Here is how I think about
it. First I imagine a large number of amino acids strung together like plastic snap-
beads, or cars of a train. (Like snap-beads or train cars, amino acids have couplings
that allow them to hook up at front and back to other amino acids, so they can form an
arbitrary sequence.) Then I imagine holding this long string of amino acids between
my two hands, tautly stretched to form a straight chain. (See Figure 27-2.) Now I let
the

        FIGURE 27-2. If you stretch a protein straight (a), and then let it go, it will
snap right back into its natural curled-up form (b), exhibiting its characteristic
tertiary structure. [Drawing by David Moser. ]




The Genetic Code: Arbitrary?                                                        675

string go. Sproingl ! 1'he crazy critter rapidly twists itself up into a tight little ball
about the size of a fist. Now you try it. Here. Grab the two ends inside the ball and
slowly pull them apart. The protein chain is resisting, of course, but if you are careful
not to jerk it, you can uncurl it without breaking it anywhere. Got it all straightened
out? Good. Now let it go. Sproing!! What do you know? It returns to exactly the same
curly shape: Now hand it back to me. Thank you.
         It seems that a protein likes to be curled up in its little ball. That three-
dimensional shape is called its tertiary structure. The tertiary structure of each type of
protein is unique. The sequence of amino acids totally determines the tertiary
structure; it is what makes the protein fold up every time into that same shape. The
one-dimensional sequence of amino acids is the protein's primary structure. Thus we
can say that a protein's primary structure determines its tertiary structure. (Some
proteins have a secondary structure as well, an intermediate level of coiling like that
of a telephone cord, but tertiary structure is the essence of a protein.)
         But so what if a protein always has a tertiary structure? So what if it folds up?
The answer is that this folded shape is what determines the kind of doer this protein
is. (If indeed it is a doer. Some proteins are not enzymes, but mere lumps that serve as
boring construction material; however, from here on out, we will be concerned only
with doer proteins-enzymes.) An enzyme's tertiary structure is characterized by
certain bumps and clefts, like the nose and ears of a person's face, except that
enzymes differ more radically than faces and are more convoluted. Certain parts of an
enzyme are called its active sites. They are where the enzyme fastens itself, leechlike,
to the lumps it is going to act on. Only by trying to fit various candidate lumps into its
active site does an enzyme eventually latch onto a lump of the proper type. (Think of
how the Prince found Cinderella by her slipper.) The enzyme and its substrate, the
lump, are often likened to a lock and a key. No wrong substrate will fit. (Actually,
wrong ones sometimes do fit under special circumstances, but we need not go into
that here.)
         An enzyme is very specific; it is tailor-made for a certain task and for no other.
Once an enzyme is hooked up to its substrates, it starts churning, like a laundromat
washing-machine into which the proper coins have been inserted. The enzyme may
rip parts off one substrate and attach them to another, it may bind two substrates
together-whatever its thing is, it does it. Then it lets go of the product or products,
which are now free to go off and drift about inside the cell, like patients after a
dentist's appointment. (See Figure 27-3.)
         The upshot of all this frenzied activity by billions of busy enzymes doing
violent things to chemical lumps inside cells is the creation and sustenance of a
unique living organism. These enzymes, these proteins, these coiled-up chains of
amino acids-these are what carry out the master plan stored passively in the cell's
genes, which are chains of-but wait! We are getting ahead of the story.




The Genetic Code: Arbitrary?                                                           676

FIGURE 27-3. The saga of Yin and Yang, two random molecules inside the cytoplasm
of a cell In (a), each drifts alone, unaware of the other. In (b), they approach an
enzyme. In (c), the enzyme recognizes that each of them fits one of its active sites, and
snares them. Then it performs its catalytic function, which in this case unites them
into one structure (d). Finally, in (e), the new YinYang unit goes on its merry way in
the cytoplasm. (Drawing by David Moser.)



The Genetic Code: Arbitrary?                                                         677

                                       *   *   *

        There is nothing more important to know about cells than how enzymes work.
They are what make cells run. But there is one other thing that is equally important.
That is, which enzymes are present, and how they got there. Not all cells have the
same set of enzymes, not by a long shot; that's why not all cells have the same
character. What's more, a given cell's complement of enzymes can change over time,
depending on both internal and external circumstances. Where do the enzymes come
from? Ultimately, of course, they come from the genes, which are like blueprints, but
that answer does not help at this point. What we need to discuss first is how enzymes
are actually built from raw materials, and then we will discuss where their blueprints
are stored.
        Remember that an enzyme is a protein, and a protein is a long chain of amino
acids linked end to end and curled up into a ball. You might think that since enzymes
are so good at taking things apart and putting things together, they would be the
protein-builders. However, the job is so delicate and specialized and critical that a
different kind of machine exists to do it. That little machine is the ribosome, of which
there are thousands in any living cell. A ribosome is partly composed of protein, but it
is also partly composed of nucleotides. Its exact composition does not matter to us,
though. After all, we are concerned only with the logic of the cell.
        Is there one ribosome for each kind of protein? Hardly. That would lead to a
terrible infinite regress. How would twenty different types of ribosome get
constructed? By twenty different types of "metaribosome"? But then how would the
twenty different types of metaribosome get constructed? You get the point. In reality,
a ribosome is not specific to any protein; it knows nothing about the various proteins
it builds. A ribosome is simply a general-purpose amino-acid hooker-upper. But then
somebody must tell it which amino acids to hook up, and in what sequence. But who?
For example, suppose the desired chain were lysine-leucine-glycine proline-cysteine-
histidinetryptophan. (I invented this sequence purely for its rhythmic quality. It is too
short to be a real protein. Proteins are usually many dozens of amino acids long, like
the seemingly endless freight trains that rumble across the midwestern plains and hoot
outside your motel room in Wibaux, Montana, late at night.) Who will tell the
ribosome to start with lysine and to finish with tryptophan?
        At the risk of seeming to invoke a different infinite regress, I shall now reveal
that there is another train, this one composed of the chemical units of Species I:
nucleotides. This train runs right through the middle of the ribosome, like a train
through a station. Its cars, taken in triplets, are what tell the ribosome which amino
acid goes first, second, third, and so on. This train is called messenger RNA, or
mRNA (where "RNA" stands for "ribonucleic acid"). A molecule of mRNA is a long
chain of A's, C's, G's, and U's. An mRNA chain is much, much longer than a protein
chain. It may consist




The Genetic Code: Arbitrary?                                                         678

FIGURE 27-4. A strand of messenger RNA, with some short double-helical regions
where it is linked to itself by hydrogen bonds (indicated by dotted lines). [Drawing by
David Moser. I

of thousands of nucleotides strung together like beads. (See Figure 27-4.) An mRNA
chain generally codes for several proteins and has special markers along it telling
where the stretches representing the various proteins begin and end. That is why there
are three special codons that do not stand for any amino acid. They act a little bit like
semicolons, in that they convey to the ribosome: "Cut this protein off right now; don't
add a single amino acid more!"
        We are coming to the crux of the matter. Where is the genetic code stored? I
have made it sound as if ribosomes "know" the code, but they do not. Although
ribosomes do the translating, they know neither language involved. How can this be?

                                       *   *   *


         Imagine yourself at the United Nations (see Figure 27-5a). An important
speech is about to be given by Mr. Na, the flamboyant ambassador from Nucleotidia.
A simultaneous interpreter of great skill, Meri Boso, is summoned. Unfortunately,
Ms. Boso has no knowledge of either the language the speech will be in or the
language it must be translated into. It looks bad! But at the last moment, just before
the speech begins, the members of a rescue team rush into the translating booth,
where they suspend from the ceiling a huge number of tiny flash cards. Each card has
on its front a word of Nucleotidian (curiously, all the words consist of three




The Genetic Code: Arbitrary?                                                         679

FIGURE 27-5. Two views of the activity of translation. In (a), Meri Boso in her U.N.
translating booth translates Mr. Na's speech from Nucleotidian into Aminoacidian,
using the flash cards dangling all about her. In (b), a ribosome chugs down an mRNA
strand, surrounded by tRNA molecules. As it reads each codon, it locates a nearby
tRNA molecule with the matching anticodon; from that tRNA, it strips off the amino
acid, attaching it to a growing protein that is emerging in its folded form into the
cytoplasm. [Drawings by David Moser. I

letters) and on its back the word's translation into the target language, which happens
to be Aminoacidian. Meri Boso is saved! All she must do is listen carefully to Mr. Na
and then, for each word she hears, find with lightning speed its flash card. Having
found the card, she deftly flips it around so that she can speak its Aminoacidian
translation in the nick of time into the microphone before her. Next word, please!
         It's no sweat, being a ribosome. All you need to do is find the right flash card
in a jiffy. But where are the flash cards in the cell? Even more to the point, what are
they? At this juncture, it seems that the genetic code has receded from view a little; it
has become more decentralized, harder to




The Genetic Code: Arbitrary?                                                         680

localize. Whereas at first one might have guessed that the genetic code was somehow
stored inside each ribosome in the chemical equivalent of a tablet or dictionary, now it
seems to lie in those flash cards. So if we want to determine how arbitrary the genetic
code is, we must determine whether the flash cards could be changed and, if so, how.
        The cell's flash cards are tRNA's (that is, transfer RNA's). The term suggests
that they are made out of the same stuff as mRNA is: A's, C's, G's, and U's. This is
true, except that some nucleotides are occasionally modified by enzymes, but for our
purposes we can ignore that fine detail. At birth, a tRNA is just an ordinary snippet of
RNA, but it is much shorter than an mRNA train. Also, quite unlike mRNA, which
stays long and snaky, a young tRNA folds up just like a protein, assuming a specific
tertiary structure. This is in contrast to mRNA, which merely forms rather aimless
curls over short stretches. The curling-up of mRNA is nonfunctional, whereas the
curling of tRNA is functional. (Or rather, we don't yet know much about mRNA's
curling; probably it is functional but in a more cryptic or subtle way.) All tRNA's fold
up into roughly the same shape: a chubby `L', rather like the bent arm of Mr. America.
At a more detailed level, however, the tertiary structures of tRNA's differ. In Figure
27-6, you will find a series of pictures of tRNA at various levels of abstraction.


        FIGURE 27-6. Transfer RNA, viewed at three levels of abstraction. In (a),
physically the most realistic, the three-dimensional structure as it has been revealed
by X-ray diffraction techniques. In (b), a more schematic "cloverleaf " representation,
showing the various hydrogenbonded loops, as well as the amino-acid attachment site
and the anticodon. In (c), the most schematic representation of all, a tRNA molecule
is portrayed in its barest functionality: a molecule labeled at one end by an anticodon
and potentially carrying at its other end an amino acid. [Drawings by David Moser. ]




The Genetic Code: Arbitrary?                                                        681

Once it is folded up, a tRNA acts like a flash card, in that it has an amino acid at one
end of the'L', and a codon at the other. Actually, it is not a codon but an anticodon. An
anticodon is to a codon as a photographic negative is to a positive, or an engraving to
a bas-relief. To make one from the other, you merely interchange A with U, and C
with G. (A and U are said to be complementary, as are C and G.) Therefore CUC and
GAG are each other's anticodons. To be more explicit about tRNA, one end of it
simply is an anticodon. The other end is a site where an amino acid can be attached.
And if you're wondering who does the attaching, you'll find out soon enough.

                                        *   *   *

         In a nutshell, a ribosome is a translating mechanism between the two
intracellular languages of Nucleotidian and Aminoacidian. The words of Nucleotidian
are codons; the words of Aminoacidian are amino acids. The mRNA is a long speech
whose sentences are written in Nucleotidian. The ribosome is a quick but ignorant
simultaneous interpreter who, guided by tRNA molecules, assembles proteins, which
are the word-by-word translations of the mRNA sentences into Aminoacidian. (By
"quick" I mean the following. Under normal conditions, a ribosome in a bacterial cell
can translate about twenty codons per second. In a rabbit cell, things are slower: a
little better than one codon per second. I have no idea why rabbits are so much slower
than -bacteria.)
         As is shown in outline in Figure 27-5b and in more detail in Figure 27-7, an
mRNA "speech" is constantly clicking through the ribosome, one codon at a time. On
encountering a new codon, the ribosome must seek out a matching tRNA, one whose
anticodon perfectly fits the codon. Of course a ribosome has no eyes, and cannot scan
about as Meri Boso does. It must try one tRNA after another (again, think of
Cinderella and her slipper). A mystery is how a ribosome can find a matching tRNA
so quickly. In any case, having found one and clicked its anticodon into position
against the mRNA codon, the ribosome snips off the tRNA's amino acid and snaps it
onto the growing protein chain; then it releases the "nude" tRNA, which is free to pick
up a new amino acid.
         This is a salient difference between the metaphorical flash cards and tRNA
molecules. Whereas flash cards can be used over and over again, each time a tRNA
molecule gets used, it has to be "recharged" with the right amino acid. Just where and
how does this take place? Which amino acid should it get charged with? How is this
determined? Who determines it? All of a sudden, these questions loom large, because
they have everything to do with the link between a codon and its amino acid. We shall
return to them shortly.
         It is now apparent that if the genetic code is stored anywhere, it is in a spread-
about fashion, distributed among the thousands of tRNA's floating in suspension in
the cell near the ribosomes. Could these tRNA's somehow be subverted? Could they
falsely guide the translation process? Certainly we




The Genetic Code: Arbitrary?                                                           682

FIGURE 27-7. The intracellular translation process, in more detail. Inside this
ribosome, one can see the matching-up of one tRNA's anticodon (while) with an
mRNA codon (black). The amino acid at the top of the tRNA has just been snapped
onto the growing chain of amino acids with a 'peptide bond", symbolized by the curly
link between rectangles. [Drawing by David Moser. ]

can imagine the UN rescue team rushing in with the wrong set of flash cards, hanging
them all up in Ms. Boso's booth, and her then translating Mr. Na's speech into a
completely inappropriate language. Could the analogue happen in a cell?. Could there
conceivably be produced an entire set of "bad" tRNA's: tRNA's with wrong amino
acids attached to them, tRNA's that would fool the ribosomes into manufacturing
nonsensical proteins? Who could perpetrate such a nasty practical joke?

                                      *   *   *

         Well, this is the stage I was at when I started drawing pictures on the
blackboard for my students. I drew a typical tRNA molecule and stated that at one
end-its AA end-it would attract a particular amino acid. But why should it attract the
right amino acid? Simple enough, I thought to myself. As with most chemical
affinities in the cell, the AA end of the tRNA would simply have the right shape. Each
tRNA would lure only the amino acid that (by the genetic code) corresponds to its
anticodon. My supposition was that for each anticodon, the tRNA that carried it
would be shaped differently at its AA end. And so that's what I drew on the board: a
tRNA molecule with




The Genetic Code: Arbitrary?                                                      683

FIGURE 27-8. My first guess at how tRNA molecules manage (nearly always) to have
just the right amino acid attached at their AA ends. In this appealing but simplistic
theory, which turned out to be completely wrong, the AA end and the desired amino
acid are like lock and key. Thus, only the desired amino acid will fit a given tRNA's
AA end. The truth of the matter is that the AA end of a tRNA is completely
nonspecific: it will accept any of the twenty types of amino acid. [Drawing by David
Moser. ]

a specific anticodon at one end and a specific "attractive shape" at the other end, a
shape that would presumably combine with just one kind of amino acid. (See Figure
27-8.)
        Here a good question arises. Why should each tRNA attract the right amino
acid for its anticodon, "right" being the amino acid defined by the genetic code? Why
couldn't some tRNA fold up in such a way as to attract some other amino acid? Or is
there some intrinsic connection between the two ends of the tRNA? Does the
anticodon, for instance, somehow "tell" the other end of the tRNA how to fold up?
This was one thought Vahe had, and it would imply that codons and amino acids
really had some direct chemical associations. But I didn't believe that for a moment.




The Genetic Code: Arbitrary?                                                     684

I told my class that neither end of the tRNA could possibly know anything about the
other. I insisted that you could surgically replace the anticodon with some other
anticodon, and the AA end would not know the difference. Conversely, you could
surgically lop off the specially shaped AA tip of the tRNA and graft on an alien AA
tip, which would then lure the wrong amino acid, thereby making the tRNA embody a
false piece of genetic code. I concluded by saying: "Since the two ends of any tRNA
are independent, the genetic code can in principle be subverted and is therefore
arbitrary. " Then I blew the chalk dust off my hands and turned to another topic.
        Well, it turns out that this picture I had drawn was right in spirit but wrong in
detail. Contrary to my supposition, all tRNA molecules have at their AA tip precisely
the same structure! For instance, the last three nucleotides at the AA tip are always
CCA (glance back at Figure 27-6). Thus, the site where the amino acid gets attached
is completely nonspecific. There is no special chemical affinity between the AA tip of
a tRNA and the amino acid that goes there! When I first found this out (after class
was over), I was somewhat at a loss. How, I wondered, does the tRNA always end up
with the right amino acid attached to it? What lures it there? Could it be the
anticodon, even though it is at the other end of the tRNA? And if so, does that mean
that there is, as Vahe surmised, some special and intrinsic link between the anticodon
and its amino acid partner? Is the genetic code, after all, inevitable?

                                       *   *   *

        By talking with biologist friends and looking in books, I found the answer. In
the end, it seemed to come out supporting my side, but matters turned out to be far
subtler and murkier than I had suspected. Although the AA end of a tRNA molecule is
indifferent to the amino acid that docks there, so that in principle it can accept any
amino acid, under normal circumstances only one amino acid will get attached.
However, this is due not to the anticodon but to the tertiary structure of another region
of the tRNA: its DHU loop. ("DHU" stands for "dihydrouridine", in case you were
curious.) This is a loop that every tRNA molecule has, and it bends around io a
characteristic shape in each different kind of tRNA. It is therefore a kind of three-
dimensional signature by which the tRNA's type can be recognized from the outside.
(Actually, as it turns out, probably considerably more is involved in tRNA recognition
than just the DHU loop, but for simplicity's sake, I will here continue to speak as if
that were the entire story.)
        But who could accomplish such recognition? Why, an. enzyme, of course -in
fact, an aminoacyl-tRNA synthetase. (Sorry about that! But despite the strangeness of
this name, you should try to remember it, because these molecules turn out to play, if
not the starring roles in our saga, then certainly pivotal roles.) Such an enzyme has
two active sites. One of them recognizes the tRNA's three-dimensional signature, and
the other looks for an amino acid. That site, unlike the AA end of the tRNA, is not
indifferent to the amino acid. It will bind one and only one amino acid-namely, the
one coded for




The Genetic Code: Arbitrary?                                                         685

by the tRNA's anticodon. To be sure, the synthetase itself never looks at the
anticodon. All it does is "sniff" the DHU loops of various tRNA's (and perhaps other
substructures as well-as I said, this is still not entirely clear) and when it finds one it
"likes", it fastens its amino acid tightly to the tRNA and bids it farewell. For each type
of amino acid, there is at least one type of synthetase.
        So here we have a funny thing. There are molecules floating around in the cell
whose purpose it is to "instruct" the tRNA's in the genetic code. They load up each
tRNA with an appropriate amino-acidic burden and then let it trudge off to encounter
a ribosome somewhere. So ... Do the tRNA's know the genetic code? No; they have to
be instructed. And who instructs them? The synthetases. Well, then ... Do the
synthetases know the genetic code? No; they merely match up DHU loops of various
shapes with amino acids. So in the end, we find out that nobody in the cell knows the
genetic code!
        Of course, that has to be an exaggeration. The truth, again, is simply that
"knowledge" of the genetic code is extremely spread out. It is shared by the entire set
of tRNA's and synthetases, and cannot be claimed by either one alone. And yet, there
is one place where one might contend that the genetic code is stored all in one piece . .
. "And where, pray tell, is that ?" you ask. Ah-it is the DNA. You might have been
wondering when we would come to DNA, usually the star in tales of molecular
biology. Well, this is the moment.

                                        *   *   *

        One might regard DNA as a big, fat, aristocratic, lazy, cigar-smoking slob of a
molecule. It never does anything. It is the ultimate "lump" of the cell. It merely issues
orders, never condescending to do anything itself, quite like a queen bee. How did it
get such a cushy position? By ensuring the production of certain enzymes, which do
all the dirty work for it. And how can it make sure that these desirable enzymes will
get produced? Ah, that is the trick.
        DNA is a set of blueprints for all kinds of cellular constituents, lumps and
doers alike. If you want to know where something in a cell comes from, the chances
are the answer is: It is coded for in the DNA. The piece of DNA that codes for some
specific entity is that entity's gene. The entity may be a protein, it may be a tRNA
molecule, or it may be some RNA that will eventually become part of a ribosome.
Whatever the constituent is, there is a gene for it in the long, twisty DNA molecule.
Indeed, that is why DNA is so long. The length of the DNA for a mere bacterium can
be a million nucleotides-and for a human being, thousands of times longer than that!
A DNA strand therefore consists of a sequence of thousands, millions, or even
billions of codons, constituting anywhere from a handful of genes to many thousands
of them, arranged sequentially, like sentences following one another in a book, or
songs in the grooves of a record.




The Genetic Code: Arbitrary?                                                           686

FIGURE 27-9. DNA at two levels of abstraction.
         In (a), an architectonic view, showing the famous double helix. The outer
winding staircases are formed of non-informational matter (sugars and phosphates),
while the inner core, represented here by hydrogen-bonded spheres, is where all the
genes are stored, defining the entire nature of the cell or organism within which all
this is occurring.
         In (b), the helices are uncoiled but no bonds are broken. This 'flattened" DNA
is then spread out like a rug on a floor, allowing you to see exactly how the bases join
up with each other in complementary pairs ("Watson-Crick bonding"). [Drawings by
David Moser. ]




The Genetic Code: Arbitrary?                                                        687

         DNA, like RNA, is made up of nucleotides, but instead of U it uses T (which
stands for "thymine"). In DNA, A and T, like C and G, are complementary. For every
strand of DNA there is a complementary strand that twists around it, making the
entire supermolecule look like a double vine (see Figure 27-9). The reason DNA does
this while RNA does not is that A and U do not fit together as tightly as A and T do,
and so the twists of any would-be RNA double helix are not as stable as those in
DNA. Actually, RNA can form a double helix for short stretches, but not for long
ones. That is also why tRNA's have short double-helical hairpin turns but are not
double helices all the way.
         I've said several times that an entity's gene is a coded version of the entity.
Now where there is code, there must be decoding. But attention: There are two
possible depths of decoding, for a stretch of DNA. (See Figure 27-10.) First of all,
you can decode it into RNA. This is the shallower way, and it is done merely by
complementation: A codes for U, T for A, and C and G for each other. Thus, DNA
stretch "TCAT" becomes RNA stretch "AGUA", which can come in handy if ever
you're parched in Paraguay. The deeper way of decoding DNA involves a second
layer of decoding (shades of the two layers of decoding of Enigma messages,
described in Chapter 21)-that is, one must go further, and decode the message
contained in the RNA. That, of course, is the job of all the Meri Bosos and their tRNA
flash cards.
         If the cell wants to make, say, a tRNA molecule, it uses only the shallow
decoding, called transcription. It finds the tRNA's gene and in effect asks the DNA-
decoding enzyme known as RNA polymerase to manufacture the corresponding
(complementary) segment of RNA. If, however, the cell wants to make a protein, it
uses both stages. First, as before, the cell gets an RNA polymerase to transcribe the
gene for the protein. The result is a long string of messenger RNA. Then this mRNA
encounters a ribosome, threads itself through it like tape through a playing head, and
as the ribosome then begins clicking down the mRNA, codon after codon, out comes
the desired protein. This second stage is called, naturally enough, translation, and
truly it lies at the dead center of life. (See Figure 27-10.)
         And so you see that in our UN scenario, Mr. Na (in whose name I hope you
recognized "mRNA") did not write his own speech. Being merely the ambassador
from Nucleotidia, he got his speech from the big boss back at home: the DNA. Mr. Na
is merely a mouthpiece, a tool. He is just reading a transcript of the DNA's speech,
and that transcript is slavishly translated by Meri Boso (whose name is "ribosome"
rotated) into Aminoacidian, according to the genetic-code flash cards. And even those
flash cards were dictated by the big boss! You see what I mean about the lazy DNA
being in control of everything?
         The DNA incorporates coded versions of'all the tRNA molecules, of all the
synthetases and polymerases, not to mention the constituents of the ribosomes. Thus
the DNA contains coded versions of its own decoders! By decoding their own genes,
the decoders manufacture more copies of themselves. You




The Genetic Code: Arbitrary?                                                        688

FIGURE 27-10. The so-called Central Dogma of molecular biology: "From DNA to
RNA to protein. " The first conversion is transcription (black arrow); the second
conversion is translation (white arrow). It is now known that reverse transcription
(from RNA to DNA) takes place in certain organisms and viruses, but reverse
translation (from proteins to either DNA or RNA) has never been observed. It is safe
to say that if it were observed, a tidal wave would be unleashed in biology, wreaking
havoc with all sorts of fundamental tenets of the science. It would entail a full-scale
return to now-discredited Lamarckian ideas about evolution. In my estimation, a
comparable event in physics would be the discovery of a method of accelerating
electrons to superluminal speeds, or perhaps the invention of a perpetual-motion
machine. [Drawing by David Moser. ]




The Genetic Code: Arbitrary?                                                       689

can see that this is a Grand Loop indeed. The genetic code is locked in, because the
decoders cannot help but produce more copies of themselves. Not only that, they also
produce enzymes that will replicate the DNA they came from, ensuring that new cells
will have exactly the same DNA, and will use exactly the same code.

                                        *   *   *

        Now, a code that is locked in is not the same thing as a code that is inevitable.
For instance, the French language is locked into France: not only do the adults in
France speak French among themselves, but also they teach it to their children.
Moreover, they publish dictionaries and grammars that stabilize the language. This
does not mean, though, that French is the only possible language in the world! The
French word for something is not intrinsically tied to that thing by some God-given
rule. French is an arbitrary code, as arbitrary as any other human language is-nest-ce
pas?
        Could it be that the genetic code is likewise changeable, despite being locked
in? How could one conceivably subvert it? What would be the chemical equivalent of
tampering with the letters in the table of the genetic code? What kind of magic wand
would I have to wave over a strand of DNA if I wanted to institute my own personal
genetic code?
        Let us set up the following hypothetical scenario. We take an ordinary
functioning cell, reach into it, and magically remove all its mRNA and tRNA. We
throw those molecules into the garbage can. Then we reach back in and remove all-
the DNA (but we hold onto it), leaving behind a lot of random flotsam and random
jetsam, including some ribosomes and some enzymes (RNA polymerases
particularly). Now these enzymes and ribosomes have nothing to do, since there is
nothing left for them to transcribe or translate. But if they will just be patient, we will
be right back! We go off and tamper with the DNA we extracted, and then we inject it
back into the unsuspecting cell. Is it possible that not only this cell, but also its
progeny, will now and forever use our new genetic code? What kinds of changes
would we have to have wrought on that piece of DNA for the cell still to function
exactly as before, except with the new code?
        What does "function exactly as before" mean, in this strange context? It means
that the cell's behavior should look, to an outside observer, as it did before. What
determines its overall functioning from that global point of view? The answer is: its
complement of proteins. Proteins, as I said early on, are what endow a cell with its
character, its personality. Given this fact, how can we ensure that our cell's external
personality is unchanged even though its internal language has been subverted?
Well, the instant we insert the altered DNA into the cell, many RNA polymerases will
start working on it. They will transcribe it into strands of RNA-both short tRNA
snippets and long mRNA trains. The tRNA snippets will fold up into their
characteristic 'L' shape. At this point, various




The Genetic Code: Arbitrary?                                                           690

synthetases encountering the fresh tRNA's will slap amino acids onto them: Then the
ribosomes will obediently use the charged-up tRNA's to translate the mRNA's. So, if
we are to produce the same proteins as before, we have to make sure of two things:

       (1) that the new tRNA's embody the new genetic code, and
       (2) that the new protein genes are written in the new code.

        Goal 1 is tantamount to making sure each tRNA has the right anticodon,
according to the new code. To achieve this goal, we just need to change three
nucleotides in each tRNA gene in the DNA. When those genes are transcribed, they
will make tRNA's that are "wrong" in just that one spot. Therefore, the first set of
changes to be made is as follows: Find all tRNA genes. In each gene, alter that DNA
codon which dictates the tRNA anticodon. To accomplish goal 2, we simply rewrite
all the "literature"that is, all the genes that code for proteins-in the new language.
(We're making "wrong" flash cards, but we're also changing Mr. Na's speech, so that
what Meri Boso says will in the end be exactly the same.)

                                       *   *   *

        Now will things really work as we hoped? For instance, will the synthetases
really do the right thing? Well, each tRNA will fold up exactly as before. (Remember
that the anticodon has no effect on the way the rest of the tRNA folds up-in particular,
no effect on the DHU loop.) Now a synthetase comes along and encounters a familiar-
seeming DHU loop. It sticks on the very same amino acid it would have stuck on
before. The enzyme has been "fooled" in exactly the way we wished. It has even
become an accomplice to our deviltry, because according to the old genetic code, the
tRNA is carrying the wrong anticodon for that amino acid, but according to the new
code, it is carrying the right one.
        If you think this scheme through, you'll see that it really works. A piece of
"alien" DNA can be inserted into a cell whose shallow decoding machinery (a set of
RNA polymerases) is present, and the cell will proceed to manufacture new deep
decoding machinery (ribosomes and tRNA's), and therewith to produce all the
proteins coded for in the alien way in the alien DNA. Collectively, these proteins will
then imbue the cell with the same external character as it had before, when it was
using the ordinary genetic code.
        Thus we have succeeded in our aim of showing that the genetic code is just as
arbitrary as Godel numbering is. And along the way, we have accomplished an
unforeseen goal: We have spelled out, in quite some detail, just what "arbitrary"
means, in the context of cellular codes.
        Gratifyingly, it turns out that inside the mitochondria of many organisms, just
such a code switch as I described above has taken place! (Mitochondria




The Genetic Code: Arbitrary?                                                        691

are semi-autonomous organelles inside cells, and their purpose is to carry out
respiration and to produce energy-carrying ATP (adenosine triphosphate) for
consumption by the host cell. They contain their own private stock of DNA, tRNA's,
ribosomes, and so forth.) The genetic code of mitochondria is very similar to the
standard genetic code: it differs in only four codons. Thus it is a dialect of
Nucleotidian, rather than a completely new language, somewhat as Joual, spoken in
parts of Quebec, is a dialect of French. Joual is just as locked in to those areas as
Parisian is to Paris. Mitochondria have their own tRNA's and their own genes, which
would not be properly understood in the main body of the cell-and yet they get along
perfectly well, thus confirming my original contention.

                                        *   *   *

         This excursion through the workings of the cell has provided only the barest
glimpse of the complexities and subtleties of the interlocking mechanisms that add up
to life. Why do there have to be so many stages and so many intermediaries? Why are
things accomplished so indirectly?
         I am reminded of a visit I made to the R. R. Donnelley plant in Chicago, where
Scientific American is printed. I was astonished by the degree of indirectness-that is,
the layers and layers of intermediary elements-of the complex machines. I kept asking
my guide about various wheels and gears and pulley systems that I saw: "What is this
for?" It always turned out that it gave the plant an extra degree of flexibility in some
way that might not have been anticipated at first. In the development of almost any
machine, the earliest model is crude. Only the most straightforward applications and
circumstances of use are taken into consideration. Then refinements introduced over
the years result in levels of complexity that make the evolved system hard, for
someone not familiar with it, to understand all at once. In fact, it may become almost
impenetrable. This certainly holds for the mechanisms of cars, airplanes, radios,
televisions, computers-even of pianos! And on a more intangible plane, this of course
applies to human language and culture, and computer software.
         In this light, it is not surprising that the cell has so many delicately balanced
mechanisms, some of which are there just to compensate for errors made by others.
Sometimes biologists and biochemists write about these things in a way that makes it
seem they have a wonderful view of the trees but have forgotten about the forest. The
way I see the machinery of the living cell, the type of counterfactual thought
experiment that comes to my mind-in *short, the view presented here-is assuredly not
the way a specialist sees it. My counterfactual thought experiment is perhaps not
experimentally feasible, but it serves the purpose of casting in stark relief the
processes that "pull" a cell's dynamic, sparking unity out of its silent, inert DNA. It
highlights the vital intermediary roles that ribosomes and tRNA's play. It serves to
focus attention on the logical crux of a cell, as we understand it: the mechanisms of
gene expression.




The Genetic Code: Arbitrary?                                                          692

        What I get out of a lucid and thorough treatise such as Lehninger's
Biochemistry is a silhouette: the shadow projected by cellular processes into the space
of information-processing concepts. I hope this is not an invalid way of looking at
things, because to me that shadow has an eerie but beautiful shape.

Post Scriptum.
         Picture this: One fine Easter morn, A-ooga and Duhhh, two strong young
protohumans of the year 198,016 B.C., are out looking for brightly colored bison
eggs. Unbeknowst to them, a ferocious green snaarfbeest, hidden in the limbs of a
nearby billaboo tree, is greedily eying them and looking forward to a couple of nice
juicy raw protohuman-burgers (without the bun). The unsuspecting pair are
approaching the tree, and as the snaarfbeest tenses up and prepares to leap upon them,
A-ooga spies it and begins to yell, "Watch out for antidisestablishmentarianism
snaarfbeest!!"-but before she has gotten all the way through that slightly awkward
definite article, the beest has leapt, and ...
         Well, I shall spare you the gory details, but one thing I can tell you for sure:
Our- two brave proto-language-users have become mere sidelines on the evolutionary
tree, and aside from the fact that their tragic tale will be retold some 200,000 years
later in a postscript concerning ergonomics and the evolution of the genetic code, they
are ciphers in the vast scheme of things. What's more, the snaarfbeest has unwittingly
done a great favor for users of rival proto-languages: It has reduced by two the
number of speakers of A-ooga and Duhhh's proto-language, and thereby strengthened
the relative position of all rival proto-languages. With enough such events, it may turn
out that the chief rival proto-language, which uses "the" in place of
"antidisestablishmentarianism" (and vice versa), will move ahead in the "Top 40"
charts for proto-languages.
         Obviously this is a ridiculous tale, but I think it gets across an idea: Efficiency
in communication really matters. Let us think about what this allegory is saying. A
language in which "the" and "antidisestablishmentarianism" were reversed might well
survive in the world, if there were no rival languages competing with it. (I can
imagine robots using it very happily.) There is nothing intrinsically wrong with "the"
being an obscure noun, and "antidisestablishmentarianism" being the most common
word in the language. Words and things don't really have intrinsic affinities for each
other. On the other hand, if you make too many of a language's high-frequency words
be long and awkward, sooner or later you are going to cross a threshold where users
of the language will be unable to keep up with the speed of events in the world, and
their survival will be imperiled. In that sense, things and their names are not totally
independent, either




The Genetic Code: Arbitrary?                                                           693

at least not if the names are to be parts of a communication system helping beings to
survive. A more compact, more elegant, more efficient language will be more able to
keep up with real-time needs. Thus the first moral is: Efficiency matters.
        A second moral, more implicit, is: Having variants matters. If you have a
number of variants, they can all fight it out and the best will survive while the weaker
ones will be pruned. Then the best ones will sprout new variants, and the same
process of selection will take place. The ratchet of evolution will advance you towards
ever more efficient variants. If, however, there is no mechanism for producing
variants, then the unique candidate will live or die simply on the basis of its own
qualities vis-a-vis the rest of the world.
        None of this is new to anyone who knows the first thing about evolution, but it
is more general than one might think. In particular, it applies directly to the question
of the uniqueness or arbitrariness of the genetic code. Various rival codes certainly
would have different efficiencies. If we presume that over many millions of years, a
bunch of rival codes arose and competed, and that out of that struggle there emerged a
winner, then it is fair to say that the winning code was not arbitrary, and that the
connections it established between codons and amino acids are preferable to a vast
number of other possibilities.

                                       *   *   *

        The allegory indicates one type of pressure: Important words should be short.
There is a pressure in all living languages toward short words (think how "car" has
supplanted "automobile", and think also of the immense number of abbreviations and
acronyms we use). There is a counterpressure, this one towards clarity and a bit of
redundancy, so that not every tiny sound is crucial. A difference between classical
Chinese and contemporary Chinese is that many words that formerly were just one
syllable long now are two. Why this shift away from shortness? Because no language
can afford to become too dense. There's not enough room in logological space-or
more precisely, in phonological space. We simply can't efficiently distinguish
between sounds that are too close together. You can't pack more than so many
monosyllabic words into a language before communication begins to suffer.
Therefore, among languages there may be fluctuations from denser to lighter, but all
languages hover about a norm, which is why translated versions of a given passage
printed side by side are all about equally long. A third pressure is of course towards
making crucial differences very salient. How risky it would be if the words for "yes"
and "no" were as close as, say, "yes" and "yef". In fact, it is rather strange that "can"
and "can't" are so close, and it is quite fascinating to observe how many phonetic
tricks we native speakers of English unconsciously employ to get around that
problem, such as glottal stops, subtle distinctions between the `a' sounds, and different
intonation patterns in order to convey whether we are being positive or negative.




The Genetic Code: Arbitrary?                                                         694

         I have been a little glib in using the word "pressure" here. What does "pressure
toward shortness", for example, actually mean? Does it mean that each speaker
actually feels an obligation to invent shorter words? No, of course not-but if by
chance a shorter way to say something comes that speaker's way, it will not be
surprising if it gets picked up and adopted, perhaps even without any conscious
awareness on the speaker's part. But this means that there must be some source of
variety. Otherwise, to speak of "pressure" has no meaning. "Pressure toward X" really
is a shorthand way of saying that variants with X will do better than ones without it.
         I must slightly qualify this remark about the meaning of "pressure". In a case
where beings with goals are capable of deliberately tailoring their behaviors,
"pressure toward X" may mean that such a being will sense that having quality X
would be more advantageous and will think up a way of getting quality X in its
behavior. Thus the source of variety in the being's behavior is internal to the being,
and a given variant is purposefully chosen by that being. Another way of looking at
this is to say that for sufficiently intelligent beings, variant possibilities can compete
in their minds, and the outcome of that simulation can determine their behavior. That
way, instead of the beings having to gamble with their lives, they just spend a little
time "programming" an internal simulation, and tailor their actual behavior according
to the results. Thus pressures are actually experienced as such, by individual
intelligent beings.
         But in most of evolution, the beings are not bright enough to be able to sense
pressures, model them internally, and respond to them consciously. Such beings
simply must accept the hand the world deals them and do the best they can. In such
cases, the meaning of "pressure" is the one involving differential selection rates in the
real world (not in a mental one) among variants with and without X-variants whose
source is external chance rather than internal reflection. One perceives the effects of
such pressure not in an individual, but in the shifting statistical makeup of a
population.
         Such is of course the case with tiny bacteria and viruses-the most primitive life
forms that presumably tried out variant versions of the genetic code way back when.
Although all of these variant codes were "viable", nonetheless some turned out to be
"more viable". Now, my column was really about the "although" clause of the
preceding sentence. All I was trying to say is that in principle, one could associate
"AGA" with any amino acid, not just arginine. The "nonetheless" clause of the
sentence includes ergonomic considerations: those that have to do with efficiency and
waste of effort. When you take ergonomics into account as well, then you realize that,
for various reasons, certain codings simply are more efficient in terms of information
theory, and those are the ones that will eventually emerge on top of the heap.

                                        *       *      *

        Along these lines, I had a very interesting letter from Robert J. Gailer of
Seattle, who pointed out to me that, contrary to what I had claimed about




The Genetic Code: Arbitrary?                                                          695

the telephone system, there was a definite ergonomic basis for the way area codes
were assigned to various cities and larger regions. The reason that New York City's
area code is 212 and San Diego's is 619, rather than the reverse, is that on dial phones
(which were universal when area codes were invented), it takes much longer to dial
619 than to dial 212. And from the point of view of long-term efficiency, when you
consider that billions of long-distance phone calls will be made, these choices are
pretty sensible. The largest cities have the shortest-dialing-time area codes. Moreover,
as Gaiter writes:

   Similar codes were assigned to non-contiguous areas to minimize confusion. It
   was theorized that if New York City had 212 while New Jersey had 213, people
   could easily get the two mixed up. By assigning 213 to Los Angeles, AT\&T
   hoped to minimize this potential source of confusion.

Thus the next two metro areas in terms of phone population, Chicago and Los
Angeles, got 312 and 213, respectively.
        Gaiter's point (similar to my point about "yes" and "yef") is that you want to
make sure that easily confusable but critically different meanings have very different
codes. Clearly, if mistakes are inevitable but some are fatal and some aren't, it's
obviously much smarter to engineer your code so that nonfatal mistakes will tend to
occur rather than fatal ones. In the case of the genetic code, this comes down to the
following. Amino acids tend to fall in families (hydrophilic-vs. -hydrophobic being
the most important class distinction, though there are others). If you make a totally
random "spelling error" in a protein (one wrong amino acid), it will almost always
destroy the desired function. However, if the mistaken amino acid belongs to the same
family as the replaced one, then the chance of salvaging some of the functional
behavior is much higher. Therefore any code that has similar codons coding for same
family amino acids will be highly favored over other codes. This favoring will have to
be a result, just to say it once more, of selection among rival codes, just as in the
allegory that opened this P.S.
        In sum, there are two different kinds of arbitrariness we are dealing with here.
In the column's sense of the term, since alternate genetic codes might survive for a
while in the absence of rivals, they are all viable and hence the one that wound up
inside our cells is arbitrary. In this P. S. 's sense of the term, since rivalry is a part of
the real world, and since selection will necessarily take place, the winning code is not
arbitrary, because it is the most efficient of a bunch of imaginable schemes. (You
were right, Vahe!) A number of people made this point very eloquently to me,
including Henriette and Miroslav Nadj, to whom I am indebted for the term
"ergonomics", Rosemarie Swanson, Nelson Max, and Barry Bunow. The most
complete response along these lines was provided by J. M. Labouygues of Clermont-
Ferrand, France, who sent me a series of articles by himself and colleagues, in which
they describe mathematical studies that have




The Genetic Code: Arbitrary?                                                            696

determined the properties of a maximally robust code-that is, one maximally resistant
to random mutations. Their work shows that the actual genetic code has those
properties, and it also shows how it could have evolved that way.
       As it turns out, the idea that the genetic code might be mathematically
optimized against mutations was first suggested-I learned this from Labouygues'
papers-about twenty years ago by my friend, the late Tracy Sonneborn, an outstanding
geneticist at Indiana University and a wonderfully alive, endlessly inquisitive, and
deeply warm human being. I laugh to think what Tracy would have thought of my
original proposal that the genetic code is arbitrary!

                                       *   *   *

        There is another way to look at the question of arbitrariness of the genetic
code, and that is to ask whether it is conceivable that one could actually carry out the
trick I suggested in the article: namely, instantaneously switch codes in some actual
living organism without impairing its life functions. I argued that one could do so by
rewriting the "literature" stored in its DNA (namely, the genes coding for all
proteins), as well as changing small pieces of its tRNA genes (namely, those coding
for the anticodon regions). To the zeroth order, this will succeed. That is, after the
switch, the two-step decoding of the DNA would produce all the same proteins as
before. If this were all that were needed to make the cell run exactly as it did before
the code switch, then we'd be in fine shape. But things are much more complicated
than that.
        When a cell's DNA looks radically different from how it used to look, but all
the same old enzymes are acting on it, something bad is sure to happen. A most
articulate discussion of this difficulty was supplied to me by Maurice Gueron in a
letter.

          A cell has to know how its DNA is organized: where a gene begins,
  where it ends, etc. For this, typographical signals are needed. Some determine,
  on the DNA, the loci which are recognized by RNA polymerases (the enzymes
  that create the messenger RNA). Others provide clues on the messenger for the
  ribosome machinery. A few such typographical signals are known-for instance,
  the 'Pribnow sequence' TATGTTG, which is involved in the recognition of the
  beginning of a gene. The existence of such signals means that there is in essence
  a second code carried by DNA: Besides the genetic code, which is used in the
  translation of nucleic acid into protein, there is also a typographical code.
          The typographical signals may well ruin any efforts to change the
  genetic code. Indeed, two things would be necessary for your scheme to work.
  First, the typographical signals should not be on a stretch of DNA that is also
  part of the genetic messages to be expressed as protein (the so-called
  "structural" genes); this is in order that rewriting the structural genes in the new
  code will not mess up the typographical signals. Second, there should be no
  possibility




The Genetic Code: Arbitrary?                                                         697

   that rewriting the structural genes will lead, by chance, to the appearance of
   new, spurious typographical signals.
            It appears that neither of these conditions is satisfied. Regarding the first,
   typographical sequences, presumably active ones, have been found within
   structural genes. As for the second, one cannot exclude the creation of spurious
   typographical sequences by a new genetic code, except through direct checking
   of all the structural genes.
            Lastly, let me make two remarks. First, the existence of two codes
   means that even with the existing machinery, some messenger RNA sequences
   must be forbidden, namely those that would have typographical significance.
   Second, the locking-in of the code means that the logic of the machinery is
   deeply interwoven with its hardware. One may still distinguish various logical
   `levels', but the connections between them will not let themselves be forgotten.

        Touche! This is precisely my own favorite point used back on me, and I am so
delighted by it. It is none other than the statement that in deep translation, you cannot
translate content alone-you must pay attention to the interaction between form and
content. It harks back to the idea of translating self-referential sentences from one
language to another-and what is more self-referential than the beautifully tangled
DNA-RNA proteins loop?
        Consider "This sentence is in English." Graduates of the A-ooga-Duhhh
Memorial School of Translation, if asked to translate it into Chinese, would produce a
Chinese sentence asserting of itself that it is in English. But that's nonsense! Either it
should refer to the original English sentence and say that that sentence is in English,
or it should refer to itself and say that it is in Chinese. This wishy-washy halfway stuff
won't do! And such remarks hold with a vengeance for the self-documenting
sentences that Lee Sallows struggled with so (see Chapters 2 and 3).
        If you substitute the word "pun" or "allusion via form" for "typographical
signal" in Gueron's commentary, and think of switching one natural language for
another, then you have another way of seeing what he is getting at. It is quite strange
but absolutely undeniable that DNA is full of puns and allusions via form. One
remarkable example is the overlapping-genes pun found in the DNA of the virus 4X
174. There, two completely different proteins were discovered to be coming from the
same section of DNA. Two genetic codes? No, nothing of the sort. Two different
reading frames. That is, by shifting the DNA over one notch, you get a new set of
codons. For example, what reads as ". . . -TGC-CAA-GGT-C ..." one way reads as
"...T-GCC-AAG-GTC- T-GCwhen you shift the frame to the right. And both ways
code for proteins indispensable to the virus's tiny quasi-life! Such incredible literary
creativity is something nobody could ever have anticipated. DNA is a marvel of self-
referential game-playing that I daresay has no rival in human literature. But then,
human literature hasn't had three billion years to develop!
        I couldn't agree more with Gueron's concluding point (which is why I put




The Genetic Code: Arbitrary?                                                             698

it in italics), and it puts me to shame, in a way. If anything, that is my own theme
song, and here I wrote an entire article attacking it! ': o be sure, I had intuited that
there was something of this sort going on in the cell, and I half-expected to be taken
severely to task by many molecular biologists for blatantly ignoring such form-
content (or structure-function) interactions but in fact, only M. Gueron did so, for
which I am very grateful to him.
         It just goes to show that if you dare knock existing establishments (in this case,
nature's chosen genetic code), thus setting yourself up as a proponent of a rebellious
disestablishmentarianism, you can be sure that somebody cleverer than you will come
back and blow your points out of antidisestablishmentarianism water, one by one,
with arguments inspired by a conservative but highly flexible the.




The Genetic Code: Arbitrary?                                                           699


                 Undercut, Flaunt, Pounce, and
                Mediocrity: Psychological Games
                         with Numbers
                                     August, 1982


IN the summer of 1962, Robert Boeninger and I, both young mathematics students
at Stanford, were riding in a bus somewhere in southern Germany on the way back
from a brief trip to Prague, when we got bored. Out of the blue, we invented a curious
game with numbers. Though the rules of our game were very simple, it was
nonetheless very tricky to play, for it involved trying to "psych each other out" in
devious ways. The rules we initially made up went like this: The game would consist
of ten turns. On each turn, we'd each choose a number in secret, and then we'd
compare them. One of us would choose a number-an integer-in the range 1-5, the
other one an integer in the range 2-6. Each of us would get to "keep" his own number,
that is, to add it to his score-provided they did not differ by 1. But in the case of two
successive integers, the player with the lower of the two numbers collected both of
them. So if I said 2 but Robert said 3, well then, I'd get 5 points, and poor Robert,
none. Very jolly! At least until I said 5 and Robert said 4. Then not so jolly.
         It seemed amusing to have the ranges not quite coincide, since it's hard to sort
out who really has the advantage. One's intuitive first impression might be that the 2-
6, or "larger", player has an advantage, but that is nicely counterbalanced by the fact
that if you name 6, you're running the risk of being undercut by your opponent's 5,
whereas your 6 itself can undercut nothing!. Moreover, the "small" player can always
name 1 safely, without any risk of being undercut.
         Although the asymmetry seemed charming, we soon decided that having equal
ranges-both 1-5-was probably preferable. And that was the way we played the game,
which I shall here call "Undercut". A table showing how much both of us stand to
lose or gain for each possible pair of choices is shown in Figure 28-la. Such a table is
known as a payoff matrix.




Undercut, Flaunt, Pounce, and Mediocrity                                             700

FIGURE 28-1. In (a), the payoff matrix for the game of Undercut, as Robert
Boeninger and I originally invented it. in each parenthesis pair, my payoff is on the
left and Robert's is on the right. In (b), Jon Peterson's way of looking at things. This
matrix exhibits the difference between Jon's profit and my profit-in other words, his
net gain over me-for each choice of moves we might make. Looked at this way,
Undercut is a zero-sum game.

         Competition was pretty fierce. The lovely thing about this game was how level
upon level of "outpsyching" could pile up in our minds. For instance, could "tease"
Robert by choosing 4 a few times in a row, trying to lure ,dim into naming 3, and just
at that moment plan to switch my move on him, jumping to 2 and outfoxing him. But
Robert of course would be most keenly aware of my ploy, and would have his own
way of playing naive, leading me on, making me think I could get away with such
tricks, and then pulling a higher-order one on me just when I least expected it.
         When I returned to Stanford from Europe that fall, I was eager to get a
computer to play this game. My friend Charles Brenner had recently written a
program that compiled frequencies of letters and letter groups (trigrams,




Undercut, Flaunt, Pounce, and Mediocrity                                            701

to be precise) in a piece of text in English (or any other language), and then, using a
random-number generator, produced pseudo-English output whose trigram
frequencies reproduced those of the input text sample. I had been very impressed by
the way that seemingly deep patterns of English could be so aptly captured by such an
algorithm, and I saw how this idea could be adapted to a game-playing program. In
particular, I was taken by the idea of getting a program to detect patterns in the move
sequences of its opponent, and then using them to generate predictions-in short,
having the computer itself try to "outpsych" its opponent. All the better if the
opponent program were also trying to do something similar to my program! The more
vicious, the better! I was in a combative mood, ready to take on all comers.
        I have vivid memories of standing over the loud line-printer where the output
would spill out, and watching the progress of games emerge line by line. We would
have our programs play games of several hundred turns, thus giving them a serious
test. My program often started out on a losing track, given that it had not yet
"smelled" any patterns in the opposing program's behavior. But sooner or later, there
would be a moment when it would appear to "catch the scent", and it would make a
decisive undercut or two in a row, and then I would see it start to surge forward, often
leaping quickly into the lead and wiping the opponent, but. This was a feeling of
overwhelming power, the power of insight defeating raw strength. It reminds me now
of one of my favorite book titles: Chess for Fun and Chess for Blood, by Edward
Lasker. I'm not anything as a chess player, but I love that title. It captures exactly that
subtle blend of goodwill and rivalry that one feels in a highly competitive game with
friends.
        Since then, I have realized how universal, how primitive, such a feeling is. It is
probably the most engrossing aspect of all sports, that feeling of pitting one strategy
against another and watching them fight it out. Dogs certainly seem to experience this
feeling with pleasure. When I play a chasing game with my friend Shandy the
Airedale, I detect in him a precise sense for how well I can anticipate his moves: in
his dodging tactics, he always stays one ply-one level of trickery-ahead of me.
Whenever I think I have caught on to his pattern, he somehow senses it, and just at
that moment shifts his strategy so that I wind up lunging for a dog that is not there.
        Oh, to be sure, he lets me win sometimes just to keep me interested. He even
has the instinct of teasing, dropping his prized stick right in front of me, acting
nonchalant about it and coolly tempting me to make a move for it. But he has it all
calculated out. He knows how quick I am, how quick he is, and what my patterns of
trying to fake him out are.
        What's more, Shandy often seems to come up with new ways of shifting his
strategy, so that I cannot simply catch onto the "meta-pattern" of his strategy-shifts
and thereby outwit him. No, there is something extremely cunning in his dog's mind,
and clearly the joyous exercise of that native




Undercut, Flaunt, Pounce, and Mediocrity                                               702

intelligence reflects a deeper quality of dogs and people in general, namely the
enormous evolutionary advantage that intelligence seems to confer on beings that
have it, in this dog-eat-dog, people-eat-people world.

                                        *   *   *

         But back to Undercut. One day, a math graduate student named Jon Peterson
who used to hang around the Stanford "comp center", as it was known back then,
challenged my program with his own. He said he had used game theory in his
program. I wasn't worried. But when I pitted my program against his, I soon saw that
there was good cause for worry. It's not that my program got trounced by his; rather
that it just never caught on to any patterns, and simply wound up more or less tying
him each time. This was baffling. Jon explained that he had computed appropriate
weights for each different choice, 1 through 5, weights that had nothing to do with the
opponent's strategy, but merely with the amount of payoff for each set of possibilities.
The payoff matrix he was talking about is shown in Figure 28-lb. It shows, for each
combination of numbers, how much Jon stands to gain relative to me. Notice the
antisymmetry-the fact that each number, when reflected across the diagonal of zeros,
changes sign, signaling that what is good for me is bad for him. (That is the definition
of a zero-sum game: when the two players' scores in any turn always cancel out.) And
of course the zeros down the diagonal mean that when we name identical numbers, it
does neither of us any good (other than carrying us one turn closer to the end).
         Since the game is completely symmetric for the two players, there can be no
winning strategy, for otherwise both players could use it and be guaranteed to beat
each other. Nonetheless, there is an optimal strategy, according to game theory, which
in a statistical sense will guarantee you long-term parity with your opponent. This
strategy is based on assigning statistical weights to the five numbers. To find those
weights, you have to solve five simultaneous homogeneous linear equations. Each
equation is based on making your expectation equal to zero. If Jon's weights for
playing 1, 2, 3, 4, 5 are, respectively, a, b, c, d, e, then my expectation, when I choose,
say, 3, will be - 2a + 5b - 7d + 2e (read straight off the third row of the payoff matrix).
Set this expectation to zero and you have one of the five equations. The other four
arise analogously. The system to solve is thus:

       (1)            -3b+2c+3d+4e=0
       (2)           3a -5c+2d+3e=0
       (3)          -2a+5b -7d+2e=0
       (4)          -3a-2b+7c -9e=0
       (5)          -4a-3b-2c+9d =0

        This amounts to inverting a 4 X 4 matrix. Jon had done so, and came up with
the following weights: 10, 26, 13, 16, and 1, for choosing 1, 2, 3, 4, and




Undercut, Flaunt, Pounce, and Mediocrity                                               703

5 respectively. Thus, according to game theory, an optimal player should play 5 very
seldom: one time out of 66. And 2 should be the most common choice. However, it
would do little good to play ten l's in a row, followed by 26 2's, 13 3's, and so on. One
must choose completely randomly, given these weights.
        Imagine a 66-sided die on ten of whose faces the number 1 appears, only one
face having 5 written on it, and so on. Each move, you must throw such a die (or a
suitable computational simulation thereof). In other words, you must avoid any and all
patterned behavior when you play according to this strategy. No matter how tempting
it might be, you must not yield! Even if your opponent plays 5 a dozen times in a row,
you must totally ignore it and merely keep on throwing your 66-sided die obliviously.
That's the way Jon's program played, and it's why my program found nothing to pick
up on. Had Jon's program ever given in to temptation and tried to outguess me, my
program would likely have picked up some pattern and twisted it back to work against
his. But his program knew nothing of temptations or teasing. It just played blindly on,
and the longer the game, the more surely it would break even. If it won, so much the
better, but it had only a fifty-fifty chance of that. That's the "optimum strategy" for
you!
        It was a humiliating and infuriating experience for me to watch my program,
with all its "intelligence", struggle in vain to overcome the blind randomness of Jon's
program. But there was no way out. I was most disappointed to learn that, in some
sense, the "most intelligent" strategy of all not only was dumb-it even paid no
attention whatsoever to the enemy's moves! Something about this seemed directly
opposite to the original aim of Undercut, which was to have players trying to psych
each other out to ever deeper levels.

                                       *   *   *

        When I saw the game so completely demolished by game theory, I abandoned
it. Recently, however, I have returned to thinking about such games in which patterns
in one's play can be taken advantage of, even if game theory in some theoretical sense
can find the optimal strategy. There is still something curiously compelling and
fascinating about the teasing and flirting and other ploys that arise in these games,
something that vividly recalls strategies in evolution, and even seems relevant to
many political situations today.
        Furthermore, there is something strikingly academic and bookish to adopting a
purely game-theoretic strategy when playing against a human opponent, especially in
the face of "teasing" strategies. Obviously, humans have more complex goals in life
than merely winning the game, and this fact determines a lot about how they play a
game. Impatience and audacity, for instance, are both important psychological
elements in human game-playing, and an optimal strategy in the ordinary game-
theoretic sense does not take those into account. Therefore I feel games of this sort are
still




Undercut, Flaunt, Pounce, and Mediocrity                                             704

important models of how people and larger organizations tackle complex challenges
and threats.
        Let me describe, therefore, some more recent variations on Undercut that I
have been experimenting with. They all involve extending the degree to which one
can go out on a limb by "baiting" one's opponent. My purpose was to encourage
teasing, which means that one player flaunts a pattern for a while, implicitly saying, "I
dare you just try an undercut!" So to $W' encourage this kind of pattern-flaunting, it
seemed reasonable to award patterns points whenever they are not picked up on by the
enemy. Let's call this version "Flaunt".
        Suppose that you and I are playing Flaunt. I say 4, you say 1. As in Undercut, I
get 4 points, you get 1. Now on my next turn, suppose I again say 4, and you say 2. If
we were playing Undercut, I would again get 4. But in Flaunt, repetitions are
rewarded. Therefore, I am given the product of my two numbers: 4 X 4, or 16. Now
suppose that on my next turn I again play 4, while you again play 2. My bravado now
earns me 4 X 4 X 4, or 64, points, while you get 2 X 2, or 4. So in these three turns, I
have gained 4+ 16+64, or 84 points, to your 1 +2+4, or 7. Of course, you have not
been oblivious to my prancing-about in front of you-you have merely been biding
your time. Now you make your move-a 3-hoping to undercut me. Too badI chose 2! I
get 5 points, you get nothing. Sorry, sucker.

       But what if I had been so dumb as to let you catch me at this? If I had indeed
said 4 this time, I would have been hoping for 256 points. But as you successfully
undercut my pattern with your 3, you get a high reward for this, namely 259 points
(your 3 points plus "my" 256 points).

                                       *   *   *

        Now Flaunt, like haircuts, can come in various styles. The one I have just
presented is the simplest. But more complex patterned behavior can also be rewarded,
if you like. I am not sure of the best way to do this, so what follows -the game I call
"Superflaunt"-is only one possible way to reward pattern-flaunting. Suppose that
instead of playing 1-2-2 against my 4-4-4 moves, you had played 2, then 1, then 2.
You might well have had a reason for doing so. Maybe it was the continuation of a
pattern for you and was worth your while keeping up for the moment. If your previous
four moves had been 2-1-2-1, then your recent three moves would have continued that
pattern. Depending on how it's scored, extending your own established pattern might
be more worthwhile than undermining my relatively new one. If 2-1-2-1-2-1-2 is
worth the product of its elements, then that amounts to 16 points. (Actually, it's worth
16 only if it was preceded by another 2-1, but that's beside the point.) By the time
you've picked up on my 4-4 pattern, maybe it seems worth it to you to let me have my
third 4 while you name one more 2, thinking that that will lull me a bit and at that
moment, you will suddenly strike, and undercut me.




Undercut, Flaunt, Pounce, and Mediocrity                                             705

         So what constitutes a pattern in this game of Superflaunt? At the moment, I'm
inclined to limit it to one fairly simple definition, although it might be possible to
have more complex definitions. The main idea is that a pattern exists when, in a given
situation, you do what you did last time you were in "that situation". So it all hinges
on what we mean by the notion of "same situation". Let's say that you have just
played x, and are about to play y. We'll say that you are making a pattern provided
that the most recent time you played x you also followed it with y. If, for instance,
your last seven moves had been 3-4-1-5-3-4-1, then to make your pattern continue,
you must play 5, and after that, 3, 4, 1, 5, 3, 4, 1, and so on. When you first establish
the sequence 3-4-1-5, you of course get no bonus points, because until the repetitions
start, there is not any pattern. Thus only when the second 4 is played has a pattern
started, and it nets you 12 (3 X 4) points. The next patterned move, 1, nets you
another 12 points, and then saying 5 gives you 60 points (as long as it is not
undercut)! But as soon as you break the pattern, your cumulative product must start
out again from scratch.
         If you had played 3-4-1-5-3-4, and were worried about the obviousness of
playing 1 now, you might choose to play 4, which, although it breaks one pattern,
establishes another pattern (viz., 4-4). Now in ordinary Flaunt, this in itself would
already be worth 16 points, but in Superflaunt only on your next 4 would you begin to
reap the benefits of your patterned playing, since only then would you have made "the
same choice" in "the same situation" twice in a row.
         A limitation of Undercut and Flaunt is that both confine your moves to a small
range. I wanted a game in which numbers of arbitrary size were permitted. It was not
too hard to come up with the following-game, which I call "Underwhelm". You and I
both think of positive integers. Now, if they are unequal and do not differ by 1, then
whoever named the lower one gets that number (the other player getting, of course,
nothing). If they differ by 1, then the namer of the upper of the two is awarded both
numbers. In that respect, Underwhelm is like a tipped-over version of Undercut
(another name for it was "Overcut"). If our two numbers are equal, then neither of us
gets anything for this turn.
         The goal can be a fixed number of points-any number. For example, 1,000
seems a good choice, although 100 or even a million will do just fine. Think about
what this does to the game. Clearly it is not useful for you to name huge numbers,
because I am likely to name a lower number and then you will get nothing while I will
get something. So there is pressure on both of us-it seems-to play fairly small
numbers. But if we stick to very small numbers, then the likelihood of being "overcut"
is fairly high. Furthermore, the scores will advance very, very slowly. If we are
progressing toward the goal of 1,000 points at a snail's pace, someone will want to
speed things up. And so someone will go out on a limb, naming a big number like 81.
Of course, doing so just once is not useful, because the other player will not have
known in advance that that 81 was coming.




Undercut, Flaunt, Pounce, and Mediocrity                                             706

        But suppose that I say 81 several times in a row. (Pattern-flaunting is not
rewarded in Underwhelm, by the way-at least not in this version.) You will soon catch
on, and may well be tempted to say 82, to overcut me. Or perhaps you will want to
make points more conservatively off my foolishness, by simply choosing numbers
close to but below 81, such as 70. Aha! Once I've managed to lure you up into my
vicinity, then of course I can start trying =to jump below you. And maybe I can even
anticipate just when you'll "bite". If so, then I can really take you to the cleaners.
        The interesting thing about Underwhelm is that by using obvious patterns as
bait to lure the opponent, either one of us can in essence establish one or more
Undercut-like games at various positions along the number line. I can set one up in
the vicinity of 81, trying to coax you into saying 82 just when I anticipate it.
Meanwhile, you may be playing a baiting game down around 30, getting 30 points
each time I extravagantly bait you with my 81, and you know that sooner or later I am
bound to try to catch you there, either going below you or overcutting you.
        What I find fascinating is how many parallel subgames of this type can arise
spontaneously in a game of Underwhelm. Particularly interesting is what happens
toward the end, when one player has a significant lead. At that point, the trailing
player will tend to play very conservatively, naming very small numbers. This means
that the possibilities for overcutting are much enhanced. Moreover, there is an entirely
psychological element to this game having to do with human impatience. Nobody
wants to dawdle to victory by choosing smallish numbers over and over again several
hundred times. Therefore, the simple quest for some variety will inevitably lead to
some quirky, daring play -every once in a while, and that will of course be
exploitable.
        Much of the spontaneous and creative teasing behavior that tends to occur in
these games has its parallels in evolution. The most picturesque and vivid portrayal
that I know of the uncanny patterns and counter patterns that are set up by living
beings competing against each other is provided by Richard Dawkins in his book The
Seth Gene. The discussion centers around the notion of an evolutionary stable
strategy, or ESS-a term due to J. Maynard Smith. An ESS is defined as: "a strategy
which, if most members of a population adopt it, cannot be bettered by an alternative
strategy". However, here, "adoption of a strategy by an individual" really means that
that individual has genes for that behavioral policy. It's not a question of choice.
        Dawkins' first example of this concept involves rival genes for two types of
aggressive behavior in a given species. The two strategies are named "hawk" and
"dove", and have the recent political connotations of those terms. If x positive points
are assigned for winning a fight, y negative points for wasting time, and z negative
points for getting injured, one can calculate,




Undercut, Flaunt, Pounce, and Mediocrity                                            707

as a function of x, y, and z, the eventual optimal balance of hawks and doves in the
population. This may be an average over time, involving swings back and forth
between mostly having hawks and mostly having doves, or it may represent an
eventual equilibrium in which the ratio is stable.
         Dawkins considers a wide variety of colorful everyday examples in human
life, carefully comparing them to strategies in the world of nonhuman evolution. Such
things as gas wars, with their price-fixing and treacherous undercutting, fall very
neatly into line with the game-theoretic analysis that he brings to bear. Some other
strategies considered are: "retaliator" (an individual who, when attacked by a hawk,
behaves like a hawk, and when attacked by a dove, behaves like a dove); "bully" (who
goes around behaving like a hawk until somebody hits back, then immediately runs
away); "proper-retaliator" (who is like a retaliator, but who occasionally tries a brief
experimental escalation of the contest). These five strategies can all be activated
simultaneously in a computer simulation of a large population, just as the strategies in
Undercut could fight each other on a computer. From such simulations, one can learn
about the optimum strategies without doing the game theory. In essence, Dawkins
maintains, this is what nature has done over eons: Vast numbers of strategies have
fought each other, nature's profligacy paying off in the long run in the development of
species with optimal strategies, in some sense of the term.
         Dawkins uses this concept to show how group selection can seem to be taking
place in a population, when in fact mere gene selection can account for what is
observed. He says:

   Maynard Smith's concept of the ESS will enable us, for the first time, to see
   clearly how a collection of independent selfish entities can come to resemble a
   single organized whole .... Selection at the low level of the single gene can give
   the impression of selection at some higher level.

        The book contains many other provocative examples of peculiar strategies that
offer sometimes frightening parallels to situations in the world of human politics,
often reminding me of the dangers of the current arms race. In fact, the connection is
made explicitly by Dawkins more than once. He refers to "evolutionary arms races"
and the survival value of deception of one species by another.
        One of the funnier parts of Dawkins' book, although it is dead serious, is
concerned with the evolution of sexuality. To show how sexuality might have
evolved, he invents "sneaky" versus "honest" gametes (fertilized eggs) and shows
how, over many generations, the former will slowly evolve into males, the latter into
females. Along the way, such amusingly named strategies are discussed as the
"domestic-bliss strategy", the "he-man strategy", the "coy" and "fast" strategies
(limited to females), and the "faithful" and "philanderer" strategies (limited to males).
Dawkins emphasizes that these are only metaphors, and are not to be taken literally
(and certainly not




Undercut, Flaunt, Pounce, and Mediocrity                                             708

anthropomorphically). When one takes them with the proper grain of salt, however,
they can enormously illuminate the mechanisms of evolution. And many of these
strategies find their counterparts in such number games as we have discussed above.

                                       *   *   *

        As I was preparing this article, I had a long phone conversation with Robert
Boeninger in which we tried out various versions of these games. One idea that
intrigued me was to play Underwhelm, but with no specific target number of points,
such as 1,000, in mind. Instead, a convention of another sort would terminate play.
My candidate for that convention was "Stop when the two players' numbers differ by
2." Thus if I say 10 and you say 8, that marks the game's end (and neither of us gets
any points on that turn).
        Robert and I tried this version out, and quickly discovered that whenever
somebody started losing, they would have no option but to go for a stalemate-a
nonterminating game. One way for the losing player to do this is to name huge
arbitrary numbers, so that they cannot be anticipated and so the condition for
termination is never met. The player who is ahead, having nothing to lose, will
cooperate by naming small numbers all the time, thereby gaining even more points
and building up even more of a lead. So you get a kind of vicious circle in which both
players wind up cooperating in a stalemate.
        Robert suggested that one way to prevent this is to add the condition that if
either player wins five turns in a row (i.e., gets a positive number of points five times
in a row), then the game is over. This prevents the player who is trailing from going
for the stalemate, because such behavior will now ensure loss. As Robert amusingly
pointed out, even if you are behind, you can start to "wind things up" by trying to win
five turns in a row, for by the time those five turns have passed, you may be in the
lead! My name for this game is "Pounce", since it made me feel like a tiger hunting
down a giraffe in the savannah, bringing down my prey in one swift, sudden move.

                                       *   *   *

        One day, several years after the Undercut episode, my sister Laura and our
friend Michael Goldhaber and I were having lunch in the Peninsula Creamery and
jotting down various trivia on napkins, as was our wont, and somehow it came to us to
play a number game involving three persons. We decided that on each turn, each of us
would choose a number in a certain range, and, since it seemed too boring to let the
biggest number win, and equally boring if the littlest number won, it became obvious
that the middlemost number should be rewarded. So we decided that on each turn,
only the "most mediocre" player's score would be allowed to increase. It




Undercut, Flaunt, Pounce, and Mediocrity                                             709

would increase, /of course, by the mediocre number; the other players' scores would
stay fixed. (A bit of a problem was posed when two players chose tying numbers, but
we found a makeshift way of handling that case.)
        Thus at the end of, say, five turns, we would all compare our scores, and the
highest one ... No, wait a minute. Why should we let the highest score win? To do that
would be, after all, contrary to the spirit of each turn. We saw quite clearly that, if the
spirit of the whole was to be consistent with the spirit of its parts, then the player
whose score was in the middle should win! We called our game "Mediocrity", but
occasionally I like to refer to it as "Hruska".
        This name was inspired by a famous remark by the then senator from
Nebraska, Roman Hruska. In those days (the early 1970's), President Nixon was
attempting to get G. Harrold Carswell appointed to the Supreme Court, against the
vehement opposition of Indiana senator Birch Bayh and others. In a radio interview
defending Carswell against his critics, Senator Hruska came out with the following
profundity:

   Even if he were mediocre, there are a lot of mediocre judges and people and
   lawyers. They are entitled to a little representation, aren't they, and a little
   chance? We can't have all Brandeises and Frankfurters and Cardozos and stuff
   like that there.

         Alas for mediocrity, Carswell's nomination was defeated. But it worked out
fine for Hruska, who shall forevermore be known as a champion of mediocrity-and
stuff like that there.
         Speaking of champs, after eating our sandwiches and drinking our thick, rich,
old-fashioned milkshakes (served in metal containers, to boot!), the three of us sat in
our booth and played a few rounds of this quirky game, and what came to our minds
but the inspired idea of determining the World Champion of Mediocrity! So we
totaled up our scores over several games, to see who had come out highest. Highest?!
Again something seemed wrong. The pervasive spirit of mediocrity that had settled on
us that day like a heavy smog urged us to deem Champion not the player who had
won the most games, not the player who had won the fewest games, but the player
who had won the middlemost number of games. Which we did, and I forget who it
was. This may be appropriate.
         At that point, a general principle seemed to be emerging, which created a
hierarchy of levels of Mediocrity. To win at Level Two (that is, our "Championship"
level), it's best to be a mediocre player at Level One (the single-game level). This
means that whereas before it was desirable to be extremely mediocre at choosing
mediocre numbers, now it's desirable to be mediocrely mediocre at choosing
mediocre numbers. How perverse! How wonderful! How wonderfully perverse! It fits
in with a general principle of perversity, a Zen-flavored principle, that applies to
many aspects of life: Try too hard, and you wind up a loser.




Undercut, Flaunt, Pounce, and Mediocrity                                               710

                                        *   * *

        After the initial session at the Peninsula Creamery in which the game of
Mediocrity was born, I worked on a number of versions of it, trying to polish it and
make it into an elegant game. I am not sure if I succeeded, but I would like to present
the rules as they presently stand.
        The major issue is how to avoid ties-not only ties at Level Zero, but at all
higher levels. My current best solution is the following: Let each player have a
slightly shifted range, relative to the other two. More concretely, let player A pick
integers from, say, 1 to 5. Then players B and C will have staggered ranges: B picks
numbers of the form n + 1/3, and C picks numbers of the form n +2/3, where n runs
from 1 to 5. Clearly, then, there can be no ties at Level Zero.
        Now what happens at Level One? Recall that a Level One game consists of
five Level Zero games, in each of which the middlemost number is awarded to the
player who chose it, with the other two players getting zero. Well, the first part of this
scoring scheme is fine, but the second part has to be modified very slightly in order to
avoid ties at higher levels. Suppose that the numbers chosen are as follows: A: 3; B: 2
3; C: 4 3. Having the middle number, A receives 3 points. B and C, however, do not
receive zero points each, but the closest positive approximation to zero that they can,
given their staggered ranges. Thus, 1/3 of a point goes to B, and 2/3 of a point to C.
        The reasoning behind this goes as follows: After five turns, each player has
received five numbers of the same form. Player A's five pure integers will add up to a
pure integer. Player' B's five numbers of the form n + 1/3 will add up to a number of
the form n +2/3, and player C's five numbers of the form n +2/3 will add up to a
number of the form n + 1/3. Thus at the next level up, B and C have exchanged roles
in terms of the form of their numbers. Consequently, the three total scores at the new
level are all of different form and cannot tie, hence there will always be a most
mediocre Level One score: a winner.
        If we now go on to consider a game at Level Two, we must award points to
each Level One game. The "winner" of a Level One game gets, of course, that
middling number of points, while the other two players once again receive the closest
positive approximations to zero possible, in their respective forms. For player A, this
means exactly zero points, as before. However, for.B it now means 2/3, and for C,
1/3. Five games at Level One constitute one game at Level Two. The heretofore tacit
"Principle of Uniformity of Levels" compels us to sum up the five Level One numbers
to produce Level Two scores. Needless to say, the same reasons as before will prevent
tie scores from arising, and so there will always be a Level Two winner.
        The same general principle will of course allow us to extend the game of
Mediocrity to any number of levels. One game of Level N+ 1 Mediocrity




Undercut, Flaunt, Pounce, and Mediocrity                                              711

consists of five Level N games. The winner of each Level N game is awarded their
Level N score, and the other two players get the minimum amount (i.e., 0, 1/3, or 2/3)
of the form of their scores at that level. These five Level N numbers are added up to
yield totals for the three players, and the middlemost one wins.
        Actually, there is nothing sacred about always having five Level N turns in a
Level N+ 1 game; the "width" could as well be four, or even two. (Multiples of three
must be avoided, since after three moves, all three players have scores that are perfect
integers, thus allowing ties.) With a width as narrow as two, this allows a very deep
(i.e., many-leveled) game to be played much more easily. For instance, with width
two, a five-level game of Mediocrity requires only 32 Level Zero turns-whereas with
the standard width of five, merely three levels of depth will require 125 Level Zero
turns. Moreover, there is nothing sacred about the Level Zero choices being confined
to numbers bounded by 5; they could run from 1 to infinity! This is just one of the
many possible variants of Mediocrity.

                                       *   *   *

        I can testify that the strategy for playing even Level Two Mediocrity gets
mighty confusing very quickly. I have played Level Three Mediocrity oil a couple of
occasions, and found it completely beyond my reach. I find this both fascinating and
frustrating. And think what it implies about world politics, if such simple games as the
ones described in this article are so baffling. How much more complex are the
"games" of international bargaining, bluffing, and war-making! All of the conceptual
messes that we have discussed above have their counterparts (only "squared", so to
speak) in international politics. As one watches these huge themes being played out
on the world stage, one can hardly help feeling like a single cell in some vast
organism whose strategy was set long ago, the consequences of which one can only
watch, hoping all will turn out for the best.

Post Scriptum.
        Suppose you are playing a very, very short game of Undercut: one turn long.
The number of points you receive will be multiplied by 1,000 and then paid to you in
dollars. What would you play? The answer must depend on what your goal is. Which
interests you more: beating your opponent, or amassing as much money as possible?
If the former is your priority, then a score of 9 to 0 (your favor) is no better than a
score of 3 to 1: Either way, you win just as fully. But if money is your desire, the
former is $6,000 more favorable than the latter. Even more striking: If you both name
5, you have




Undercut, Flaunt, Pounce, and Mediocrity                                            712

a tie game-a big disappointment for someone out to win, but for someone out for
money, $5,000 is a fine take.
        There is thus a big difference between the original payoff matrix for Undercut
and Jon Peterson's modified matrix (both shown in Figure 28-1). Jon's matrix looks at
the game solely from the point of view of someone who wants to beat the other
player. By taking, the difference between payoffs, Jon managed to convert Undercut
into a zero-sum game, which he knew to be tractable by the methods of game theory.
But if he had left it as Robert and I had formulated it to begin with, it would not have
been so easy.
        In fact, the original (non-zero-sum) formulation of Undercut subsumes the
most famous of all non-zero-sum games: the Prisoner's Dilemma (a treacherous
Gordian knot with which the next few chapters deal). I have extracted, in Figure 28-2,
just a small fragment of the Undercut payoff




FIGURE 28-2. A portion of the original Undercut payoff matrix, showing how
Undercut contains a Prisoner's Dilemma matrix. (In fact, it contains several.)

matrix, geometrically rearranged but otherwise intact. This miniature payoff matrix
has virtually all the same mathematical qualities as does the standard Prisoner's
Dilemma payoff matrix (Figure 29-1). Thus Undercut actually poses a more severe
problem than Jon Peterson said. His trick of subtracting one person's payoff from the
other's will turn any symmetric game into a zero-sum game, which is tractable by
standard techniques of game theory. But that ignores significant aspects of the
original game; in particular, for any normal person, losing by 3 to 5 (and getting
$3,000) is precisely as good as winning by 3 to 1 (also getting $3,000). But in Jon's
matrix, these two events are as opposite as night and day-as opposite as -2 and +2.

                                       *       *      *

        One real-life counterpart to Undercut is given by the following amusing
observation. The long-distance telephone rates get much cheaper at 11:00 at night,
and so as 11 approaches, the lines get less and less busy, until suddenly, when the
hour strikes, the lines get very crowded. In some parts of the country, this "rush hour"
prevents you from being able to get a line




Undercut, Flaunt, Pounce, and Mediocrity                                            713

at all, which is molt annoying. So you have the option of calling just before 11 and
getting an expensive line, or calling just after 11 and taking your chances. Maybe you
decide that it's better to call just before 11, and pay the extra amount for the security
of getting through. But if everybody thinks of this strategy, then calling just before 11
is self-defeating! So then you have to start pushing your calls back earlier, perhaps
even into a more expensive time period ... I guess this is just a new variant on the old
"Nobody ever goes there any more because it's always too crowded" joke.
        Games of this sort and jokes do indeed have a lot in common. In an article in
the British journal Manifold titled "A Pandora's Box of non-Games", Anatole Beck
and David Fowler set forth a panoply of rather silly games that are halfway between
true games and pure jokes. The tragedy is that so many of them resemble current
global political behavior. For instance, consider the game they call Finchley Central:

   Two players alternate naming the stations on the London Underground. The
   first to say 'Finchley Central' wins. It is clear that the `best' time to say 'Finchley
   Central' is exactly before your opponent does. Failing that, it is good that he
   should be considering it. You could, of course, say 'Finchley Central' on your
   second turn. In that case, your opponent puffs on his cigarette and says, 'Well....
   ' Shame on you.

        Another amusing game, quite similar to the ones described in the column, is
called Penny Pot

   Players alternate turns. At each turn, a player either adds a penny to the pot or
   takes the pot. Winning player makes first move in next game. Like Finchley
   Central, this game defies analysis. There is; of course, the stable situation in
   which each player takes the pot whenever it is not empty. This is a solution?

At the end of their article, Beck and Fowler add:

   M. Henton of New Addington noted with horror that there is an isomorphism
   between Finchley Central- and the game commonly known as 'Nuclear
   Deterrent'. 'It occurs to me that we should work very fast to analyse the non-
   games, before we are left with a non-world.'

        Several readers wrote in to tell me that they had worked out by game theory
the optimal strategy for playing my game of Underwhelm, and that they had found it
involves playing only the numbers between 1 and 5, in the ratios 25:19:27:16:14.
Numbers higher than 5 should never be played at all! This was a surprise to me,
taking away most of the interest of the game. Oh, well ... as they say in game theory,
"You win some, you lose some."




Undercut, Flaunt, Pounce, and Mediocrity                                                 714

